\chapter{Matrix Factorizations}
\label{ch:matrix-factorizations}

The third source of CY categories is the Landau--Ginzburg model: a polynomial $W \colon \mathbb{C}^n \to \mathbb{C}$ whose category $\MF(W)$ of matrix factorizations is CY of dimension $n - 2$.  For ADE singularities $W = x^N + y^2 + z^2 + w^2$, $\MF(W)$ is CY$_2$ and $\Phi$ (Theorem~CY-A$_2$) produces chiral algebras related to $\cW_N$-algebras.  The LG/CY correspondence $\MF(W) \simeq D^b(\Coh(X_W))$ provides a consistency check: $\Phi(\MF(W))$ and $\Phi(D^b(\Coh(X_W)))$ (Chapter~\ref{ch:derived-cy}) must agree.

\section{Matrix factorizations and Landau--Ginzburg models}
\label{sec:mf-lg}

\section{The CY structure on $\MF(W)$}
\label{sec:mf-cy-structure}

\section{Quantum chiral algebras from singularities}
\label{sec:mf-qca}

\section{ADE singularities and W-algebras}
\label{sec:ade-w-algebras}

The connection between ADE singularities and W-algebras provides a key test case for the CY-to-chiral functor.  For the $A_{N-1}$ singularity $W = x^N$ in one variable, $\MF(x^N)$ is CY of dimension $-1$ (formal), so the relevant setting is the two-variable singularity $W = x^N + y^2$, giving $\MF(W)$ of CY dimension $0$ (semisimple).  To reach a CY$_2$ category related to $W_N$, one considers the threefold singularity $W = x^N + y^2 + z^2 + w^2$ in four variables, yielding $\MF(W)$ of CY dimension $2$.
% RECTIFICATION-FLAG: The original claimed MF(x^N) "should produce a
% quantum chiral algebra related to W_N" without specifying the ambient
% dimension. The CY dimension of MF(W) depends on the number of variables.
% For the W-algebra connection, one needs a CY_2 category (to apply
% Theorem CY-A_2), which requires 4 variables.
